\documentclass[12pt]{article}
\usepackage[utf8]{inputenc}
\usepackage[T1]{fontenc}
\usepackage{hyperref}
\title{Sunam Language Resources for K-12 Curriculum Development}
\author{Compiled by Claude Sonnet 5.5 (claude-sonnet-5-5) for Andr\'as Kornai}
\date{}
\begin{document}
\maketitle

\section{Introduction}
Sunam (ISO 639-3 ssk) is a little-documented Sino-Tibetan language spoken in Sunnam village, Kinnaur district, Himachal Pradesh, India. Its classification is disputed. Negi (2020) treats it as a Bodic-branch Tibeto-Burman language, and Widmer (2018) places it in the Eastern branch of West Himalayish. Wikipedia's Kinnauri article lists it as a Kinnauri dialect, but that claim is marked as needing a citation. Speakers also use Shumcho, Hindi and Kinnauri, and Hindi use has grown since 1990.

\section{Primary Educational Resources}
No primers, textbooks or school materials in Sunam turned up in the search.

\section{Linguistic and Reference Materials}
\begin{itemize}
\item Negi, Language Documentation and Description 19 (2020), pp.\ 249--255. A language snapshot of Sunam: \url{https://www.lddjournal.org/article/id/1232}
\item Wikipedia, Sunam language: \url{https://en.wikipedia.org/wiki/Sunam_language}
\item The Linguistic Landscape of the Indian Himalayas (OAPEN). It covers Kinnauri and neighbouring varieties, but not Sunam itself, so it is useful for regional context only: \url{https://library.oapen.org/handle/20.500.12657/76653}
\end{itemize}

\section{Religious and Community Resources}
No religious or community resources turned up in the search.

\section{Digital and Audio Resources}
No audio, video or software resources turned up in the search.

\section{Curriculum Development Context}
Negi describes Sunam as essentially undescribed. There are no teaching materials, so a curriculum would first need basic documentation: a grammar sketch, a lexicon, an orthography and recordings. The snapshot reports that the language is still passed on to children and used in daily life, which is a favourable starting point.
\end{document}